\subsubsection{Embedding model}\label{sec:embedding-model}
The thesis embeds text using a sentence-transformers model, which is a continuation of the BERT model class. BERT is an encoder-only transformer \citep{devlin2019bert}, and it had to be adapted before it could be used to compare two sentences and measure their similarity. Three properties of the original model stand in the way.

Firstly, BERT is not trained to recognise similarity. Alongside masked token prediction, it is trained on next sentence prediction, which asks how likely it is that one sentence follows another in a document. This is a wholly different relation than the one required here. A sentence that follows another need not resemble it in any way, and semantically similar sentences do not necessarily follow each other.

Second, and even if next sentence likelihood were the relation of interest, BERT scores a pair by passing both sentences through the model together. Individual sentences therefore cannot be assigned a single reusable embedding vector, and each sentence has to be scored against every other sentence in the corpus. This results in an explosion of compute and time, which rules the approach out for a memory store that is searched on every Think turn.

Lastly, the raw encodings do not separate cleanly under cosine similarity. BERT's embeddings are anisotropic, meaning they occupy a narrow cone of the vector space rather than spreading through it \citep{ethayarajh2019contextual}, so arbitrary sentences return high similarity scores.

SBERT was created as a solution to these issues \citep{reimers2019sentencebert}. It takes the mean of all token vectors to generate a single fixed-size semantic representation of a sentence, and it fine-tunes the underlying BERT through a siamese network on sentence pairs, so that similar sentences are pulled together in the vector space and dissimilar ones pushed apart. The pooling on its own is not what makes this work: averaged embeddings from an un-tuned BERT perform worse on similarity tasks than averaged GloVe vectors \citep{reimers2019sentencebert}. It is the fine-tuning that makes the cosine similarity between two sentence vectors meaningful.

The model used in this thesis, \texttt{all-MiniLM-L6-v2}, is a six-layer reduction of the distilled MiniLM transformer \citep{wang2020minilm}, fine-tuned in the sentence-transformers framework on roughly 1.17 billion sentence pairs. It is used in two places. First, to retrieve semantically similar memories for the Think prompt based on the agent's task, and second, to assign each memory to one of the differentially private labels in the shared memory pipeline described in Section \ref{sec:differential privacy mechanism}.

The model was chosen for its cost, since the retrieval step is called once per Think turn across 812 tasks per pass-through, which means the per-call cost compounds quickly. At 22.7 million parameters and 384 dimensions it runs offline and deterministically. MTEB \citep{muennighoff2023mteb} reports that no single embedding model dominates across task types, so the choice is a trade-off rather than a ranking, and \texttt{all-MiniLM-L6-v2} sits at the small and fast end of it. While a larger encoder would likely retrieve more accurately, this thesis compares configurations against each other under one fixed retrieval setting, so holding retrieval quality constant across conditions matters more than maximising it.
